\section{Direction Selection: Why the Rotation Goes Clockwise or Counterclockwise}
\label{sec:direction}

A natural question that the previous sections left open is: what determines whether the formation
rotates clockwise ($\omega=-\sqrt{k_2}$) or counterclockwise ($\omega=+\sqrt{k_2}$)? The controller
\eqref{eq:ctrl7} is symmetric under reflection---there is no built-in preference for either
sense---yet every numerical run eventually settles into one of the two directions. This section
resolves the question using the concept of \emph{directed circulation (angular momentum)} as an
order parameter.

\subsection{Reflection symmetry: two mirror limit cycles}

The closed-loop dynamics in the target frame,
\begin{equation}
  \dot{\txi}_i=\phi_i-k_1\txi_i+\tvell_i,\qquad
  \dot{\tvell}_i=-k_2\txi_i,
  \label{eq:dirn_dyn}
\end{equation}
is invariant under the reflection $J\mapsto-J$, i.e., under the transformation
$(\txi_i,\tvell_i)\mapsto(R\txi_i,R\tvell_i)$ for any reflection matrix $R$. This follows from the
rotation covariance of the repulsive forces, $\phi_i(R\txi)=R\phi_i(\txi)$. Consequently every
rotating steady state has a mirror image with the opposite rotation sense:
\begin{equation}
  \omega=+\sqrt{k_2}\quad\text{(counterclockwise)}\qquad\Longleftrightarrow\qquad
  \omega=-\sqrt{k_2}\quad\text{(clockwise)}.
  \label{eq:mirror}
\end{equation}
These two limit cycles are distinct attractors in the phase space, and the question of direction
reduces to which basin of attraction the initial condition falls into.

\subsection{Order parameter: directed circulation}
\label{sec:orderparam}

The natural order parameter that distinguishes the two basins is the \emph{directed angular
momentum} (circulation) about the target:
\begin{equation}
  L=\sum_{i\in\NN} (x_i-x_0)\times(v_i-v_0)
    =\sum_{i\in\NN} \txi_i\times\tvell_i .
  \label{eq:angmom}
\end{equation}
On a rotating steady state $\tvell_i=\omega J\txi_i$, so
\begin{equation}
  L=\sum_{i\in\NN} \txi_i\times(\omega J\txi_i)
    =\omega\sum_{i\in\NN}\norm{\txi_i}^2,
  \qquad\text{hence}\qquad
  \operatorname{sign}(L)=\operatorname{sign}(\omega).
  \label{eq:Lsign}
\end{equation}
Thus $L$ correctly labels the two rotating steady states. It should not, however, be interpreted as
a conserved quantity: away from a steady rotation, $L(t)$ need not be constant, and the full level set
$L=0$ is not generally an invariant manifold. What is invariant is the smaller reflection-symmetric
subspace containing perfectly symmetric initial conditions, such as the regular-hexagon setup used
in the paper's typical experiments.

\subsection{Numerical evidence: direction = sign of initial circulation}
\label{sec:dirn_num}

The prediction was tested by initializing each vehicle with a small tangential perturbation
$v_i(0)=\varepsilon J\,x_i(0)$, so that $\sgn L(0)=\sgn\varepsilon$, and sweeping $\varepsilon$ over
twelve orders of magnitude in both signs; the runs are tabulated in
Section~\ref{S-sec:sim_direction} of the companion. The result is unambiguous and matches
\eqref{eq:Lsign}: every $\varepsilon\ne0$ run settles onto a rigid rotation with
$\sgn\omega=\sgn\varepsilon$ and $\abs{\omega}=\sqrt{k_2}$, down to $\varepsilon=\pm10^{-12}$.

Two conclusions follow. The \emph{sign} of the initial circulation selects the rotation sense, and its
\emph{magnitude} does not: perturbations differing by eleven orders of magnitude converge to the same
steady $\omega$ to eight digits. Any nonzero circulation is therefore amplified by the transient into
a full rigid rotation, which is what makes $L$ an order parameter rather than merely a correlate of
the final state.

\subsection{Spontaneous symmetry breaking at zero circulation}
\label{sec:ssb}

The paper's typical initial condition (regular hexagon, zero initial velocity estimates) has
$L(0)=0$ and lies in a reflection-symmetric subspace. This symmetric subspace is the degenerate case
in which neither clockwise nor counterclockwise rotation is preferred by the equations. It is not
correct to say that every state with $L=0$ remains there, nor can a nontrivial rigid rotation have
zero circulation, because \eqref{eq:Lsign} gives $L=\omega\sum_i\norm{\txi_i}^2$. In exact arithmetic
a perfectly symmetric trajectory would preserve its imposed reflection symmetry. In finite-precision
arithmetic, however, floating-point noise at the $10^{-16}$ level provides a symmetry-breaking
perturbation that is amplified over a long time scale.

The behaviour of the exactly symmetric initial condition is reported in
Section~\ref{S-sec:sim_ssb}. The outcome is more subtle than a slow drift toward one of the two mirror
attractors: within the horizon tested ($6000$\,s) the symmetric run does \emph{not} lock onto a rigid
rotation at all. The per-vehicle rates $\omega_i$ retain a large spread, so the aggregate
$\bar\omega$ is not a meaningful rotation rate for this trajectory, and no approach to
$|\omega|=\sqrt{k_2}$ is observed. What the run does exhibit is a persistent negative circulation
$L\approx-66$, i.e.\ a definite \emph{sense} of net circulation, without convergence to a rigid
orbit.

The correct reading is therefore weaker than a completed pitchfork. Floating-point noise does break
the reflection symmetry---the sign of $L$ is selected and then maintained---but the symmetric
manifold is only slowly repelling, and the trajectory lingers near it in a non-rigid, irregularly
circulating state for far longer than the transient of any $\varepsilon\ne0$ run
(which locks by $t\approx400$\,s, cf.\ Section~\ref{S-sec:sim_direction}). Whether the symmetric initial
condition eventually reaches one of the two mirror limit cycles is not settled by this experiment.

\begin{remark}[Practical consequence for the paper's figures]
The direction observed in Figs.~2--4 of \cite{kou2022cooperative} should not be read as intrinsic:
the regular-hexagon, zero-velocity initial condition is precisely the critical case, and the
selected sense of circulation is a product of floating-point noise, so a different machine,
compiler, or integrator may produce the opposite sense. Two caveats apply. First, as just noted, an
exactly symmetric run need not settle onto a rigid rotation on the time scales usually plotted, so
figures generated from it may show a slowly circulating non-rigid state rather than a limit cycle.
Second, this concerns only the \emph{direction} and the \emph{approach}: the robust
conclusions---$|\omega|=\sqrt{k_2}$ on any rigid rotation, and the local orbital stability of the
rotating branch---are unaffected, and they are confirmed for every $\varepsilon\ne0$ initial
condition tested.
\end{remark}

\subsection{Summary of the direction-selection mechanism}
\begin{enumerate}[leftmargin=2.5em]
  \item The system has reflection symmetry, so the two rotation senses are mirror-image attractors.
  \item The \emph{directed circulation} $L=\sum_i\txi_i\times\tvell_i$ is the order parameter:
        $\operatorname{sign}(L)=\operatorname{sign}(\omega)$ on any rotating steady state.
  \item In the perturbation family tested here, any nonzero initial circulation (even at
        $10^{-12}$) determines the direction uniquely.
  \item When the initial circulation is exactly zero, floating-point noise breaks the reflection
        symmetry and selects a persistent \emph{sense} of circulation, but the symmetric manifold is
        only weakly repelling: no rigid rotation is reached within the horizon tested, so the
        analogy with a completed pitchfork bifurcation is suggestive rather than established.
\end{enumerate}

